%==============================================================================
% IEEEtran LaTeX Template for ClinX Clinical Paper
%==============================================================================
\documentclass[10pt, conference, compsocconf]{IEEEtran}
% For journal: \documentclass[10pt, journal, compsoc]{IEEEtran}

\usepackage{cite}
\usepackage{amsmath,amssymb,amsfonts}
\usepackage{algorithmic}
\usepackage{graphicx}
\usepackage{textcomp}
\usepackage{xcolor}
\usepackage{booktabs}
\usepackage{multirow}
\usepackage{siunitx}
\usepackage{hyperref}
\usepackage{balance}

%==============================================================================
% MACROS & NOTATION
%==============================================================================
\newcommand{\nt}{Nucleotide Transformer}
\newcommand{\esm}{ESM-2}
\newcommand{\acmg}{ACMG/AMP}
\newcommand{\vus}{VUS}
\newcommand{\cds}{CDS}
\newcommand{\clinvar}{ClinVar}
\newcommand{\uniprot}{UniProt}
\newcommand{\grch}{GRCh38}
\newcommand{\pval}{$p$-value}
\newcommand{\etal}{\textit{et al.}}
\newcommand{\ie}{\textit{i.e.}}
\newcommand{\eg}{\textit{e.g.}}
\newcommand{\etc}{\textit{etc.}}
\newcommand{\wrt}{\textit{w.r.t.}}
\newcommand{\fig}{Fig.}
\newcommand{\tab}{Table}
\newcommand{\alg}{Algorithm}
\newcommand{\secn}{Section}
\newcommand{\eqn}{Eq.}

%==============================================================================
% DOCUMENT
%==============================================================================
\begin{document}

\title{ClinX Clinical: An Integrated Deep Learning and ACMG Framework\\for Clinical-Grade Genomic Variant Interpretation}

\author{
\IEEEauthorblockN{
First Author\IEEEauthorrefmark{1},
Second Author\IEEEauthorrefmark{2},
Third Author\IEEEauthorrefmark{1},
and Senior Author\IEEEauthorrefmark{3}
}
\IEEEauthorblockA{
\IEEEauthorrefmark{1}Department of Computer Science, Institution, City, Country\\
email@domain.com
}
\IEEEauthorblockA{
\IEEEauthorrefmark{2}Department of Genetics, Institution, City, Country\\
email@domain.com
}
\IEEEauthorrefmark{3}Corresponding Author: Senior Author, Department, Institution, City, Country\\
email@domain.com
}

\maketitle

\begin{abstract}
Clinical interpretation of genomic variants remains a critical bottleneck in precision medicine, with 40--60\% of variants classified as Variants of Uncertain Significance (\vus). Existing automated tools lack integration of state-of-the-art deep learning models with calibrated scoring and complete \acmg\ 2015 rule implementation. We present \textbf{ClinX Clinical}, an end-to-end clinical-grade variant interpretation system that integrates: (1) \nt\ for DNA-level pseudolikelihood scoring with gene-specific empirical calibration; (2) \esm\ for protein-level mutational effect prediction; (3) robust transcript mapping via Ensembl REST API with curated fallbacks for clinical genes; (4) a complete Bayesian \acmg\ classification engine implementing all 28 evidence criteria; and (5) \clinvar/\uniprot\ evidence integration. Deployed as a FastAPI/React system with Google Colab A100 GPU acceleration, ClinX Clinical achieves \textbf{XX\%} concordance with \clinvar\ expert classifications on a benchmark of 50 \textit{BRCA1}/\textit{TP53} variants, reducing \vus\ rate by \textbf{XX\%} while maintaining \textbf{>90\%} pathogenic recall. The system provides a scalable, transparent framework for clinical variant triage.
\end{abstract}

\begin{IEEEkeywords}
Genomic Variant Interpretation, ACMG/AMP Guidelines, Deep Learning, Nucleotide Transformer, ESM-2, Clinical Bioinformatics, Precision Medicine, BRCA1, TP53
\end{IEEEkeywords}

%==============================================================================
% I. INTRODUCTION
%==============================================================================
\section{Introduction}
\label{sec:intro}

% [~800 words]
% 1. Clinical variant interpretation crisis - VUS rates, manual curation bottleneck
% 2. Current landscape - CADD, REVEL, AlphaMissense, InterVar, Franklin, Varsome
% 3. Critical gaps (4-5 bullet points)
% 4. Our contributions (5 numbered contributions - C1 to C5)
% 5. Paper organization

%==============================================================================
% II. RELATED WORK
%==============================================================================
\section{Related Work}
\label{sec:related}

\subsection{Deep Learning for Variant Effect Prediction}
% Nucleotide Transformer, ESM-2, AlphaMissense
% Gap: no dual-model integration with calibration

\subsection{ACMG/AMP Automation Tools}
% Table comparing InterVar, Franklin, Varsome, ClinX
% Criteria coverage, scoring method, DL integration, deployment

\subsection{Clinical Variant Databases}
% ClinVar, UniProt, gnomAD

\subsection{Calibration in Genomic Deep Learning}
% Temperature scaling, isotonic regression
% Gap: no gene-specific empirical calibration

%==============================================================================
% III. METHODOLOGY
%==============================================================================
\section{ClinX Clinical Framework}
\label{sec:methodology}

% Fig. 1: System Architecture Pipeline

\subsection{System Architecture}
\label{sec:architecture}
% FastAPI backend, React frontend, Colab A100 GPU, <10s latency

\subsection{Transcript Mapping and CDS Localization}
\label{sec:transcript}
% Algorithm 1: Genomic-to-CDS mapping
% Table I: Curated fallback mappings (22 variants)

\subsection{DNA-Level Scoring: Nucleotide Transformer with Calibration}
\label{sec:dna}
% Algorithm 2: Calibrated DNA scoring
% Fig. 2: Calibration curves
% Table II: Gene-specific calibration parameters

\subsection{Protein-Level Scoring: ESM-2}
\label{sec:protein}
% Algorithm 3: Protein mutational effect scoring

\subsection{ACMG/AMP 2015 Bayesian Classification Engine}
\label{sec:acmg}
% Table III: 28 criteria implementation status
% Algorithm 4: Bayesian classification
% Fig. 3: ACMG decision flowchart

\subsection{Clinical Evidence Integration}
\label{sec:evidence}
% ClinVar eUtils API, UniProt REST, caching strategy

%==============================================================================
% IV. EXPERIMENTAL SETUP
%==============================================================================
\section{Experimental Setup}
\label{sec:experiments}

\subsection{Benchmark Dataset}
% Table IV: 50 variants (25 BRCA1, 25 TP53) with ClinVar truth

\subsection{Baseline Methods}
% Table: CADD, REVEL, AlphaMissense, InterVar, ClinVar Consensus

\subsection{Evaluation Metrics}
% Accuracy, Macro-F1, Pathogenic Recall, VUS Rate, Cohen's kappa, ECE

\subsection{Hardware and Reproducibility}
% Colab A100, CPU specs, latency, code availability

%==============================================================================
% V. RESULTS
%==============================================================================
\section{Results}
\label{sec:results}

\subsection{Overall Classification Performance}
% Table V: Comparison table
% Fig. 4: Confusion matrix

\subsection{Per-Class Analysis}
% Table VI: Per-class P/R/F1

\subsection{VUS Reclassification Impact}
% Fig. 5: Sankey diagram

\subsection{Ablation Study}
% Table VII: Component ablation
% Fig. 6: Ablation bar chart

\subsection{Calibration Effectiveness}
% Fig. 7: Reliability diagrams

\subsection{Computational Performance}
% Table VIII: Latency breakdown

%==============================================================================
% VI. DISCUSSION
%==============================================================================
\section{Discussion}
\label{sec:discussion}

\subsection{Clinical Utility}
% Triage, VUS resolution, transparency, scalability

\subsection{Comparison to State-of-the-Art}
% vs InterVar, Franklin, AlphaMissense

\subsection{Limitations}
% 1. Gene scope, 2. Population data, 3. Splicing, 4. Circularity, 5. Frameshift protein, 6. Clinical validation

\subsection{Ethical and Regulatory Considerations}
% CDS guidance, clinician-in-loop, privacy, bias

\subsection{Future Work}
% Multi-gene panels, gnomAD, SpliceAI, prospective trial, FHIR

%==============================================================================
% VII. CONCLUSION
%==============================================================================
\section{Conclusion}
\label{sec:conclusion}
% Summary, key results, clinical impact, availability, final statement

%==============================================================================
% REFERENCES
%==============================================================================
\bibliographystyle{IEEEtran}
\bibliography{references}

%==============================================================================
% BIOGRAPHIES (for journal)
%==============================================================================
%\begin{IEEEbiography}[{\includegraphics[width=1in,height=1.25in,clip,keepaspectratio]{author1.jpg}}]{First Author}
%Biography text...
%\end{IEEEbiography}

\end{document}